\documentclass[11pt,a4paper]{article}
%-------------------------
% Basic packages
%-------------------------
\usepackage[utf8]{inputenc}
\usepackage[T1]{fontenc}
\usepackage{lmodern}
\usepackage{microtype}
\usepackage{geometry}
\geometry{margin=1in}

% Mathematics
\usepackage{amsmath,amssymb,amsthm,mathtools,bm}
\usepackage{physics}    % convenient bra-ket, derivatives, etc.
\usepackage{siunitx}    % units
\usepackage{braket}     % alternate bra/ket macros
\usepackage{mhchem}     % chemistry notation if needed

% Figures & graphics
\usepackage{graphicx}
\usepackage{tikz}
\usepackage{pgfplots}
\pgfplotsset{compat=1.18}

% References, hyperlinks, and clever names
\usepackage[colorlinks=true,linkcolor=blue,citecolor=blue,urlcolor=blue]{hyperref}
\usepackage[nameinlink]{cleveref}

% Lists and layout
\usepackage{enumitem}
\setlist{itemsep=2pt,parsep=2pt}

% Theorem environments
\theoremstyle{plain}
\newtheorem{theorem}{Theorem}[section]
\newtheorem{lemma}[theorem]{Lemma}
\newtheorem{proposition}[theorem]{Proposition}
\newtheorem{corollary}[theorem]{Corollary}

\theoremstyle{definition}
\newtheorem{definition}[theorem]{Definition}
\newtheorem{example}[theorem]{Example}

\theoremstyle{remark}
\newtheorem{remark}[theorem]{Remark}

%-------------------------
% Useful macros
%-------------------------
% \newcommand{\dd}{\mathrm{d}}                      % differential
\newcommand{\ii}{\mathrm{i}}                      % imaginary unit
\newcommand{\eexp}{\mathrm{e}}                    % exponential base
\DeclareMathOperator{\sgn}{sgn}

% Shortcuts for common physics notation (overrides safe if physics package present)
\renewcommand{\vec}[1]{\boldsymbol{#1}}

%-------------------------
% Bibliography (biber/biblatex) - optional
%-------------------------
% \usepackage[backend=biber,style=phys]{biblatex}
% \addbibresource{references.bib}

%-------------------------
% Document metadata
%-------------------------
\title{Execise 2}
\author{Zhuge X.K.}
\date{\today}

%-------------------------
% Document
%-------------------------
\begin{document}

\maketitle
\section*{Exercise 3.2}
At the interface between the normal metal and the superconductor, Andreev reflection occur. The BdG equation can be written as:
\begin{equation*}
    \ii \hbar \partial_t 
    \begin{pmatrix}
        f\\
        g
    \end{pmatrix}
    =
    \begin{pmatrix}
        -\frac{\hbar^2}{2m}\nabla^2 - \mu(x) + V(x) & \Delta(x)\\
        \Delta^* & \frac{\hbar^2}{2m}\nabla^2 + \mu(x) - V(x)
    \end{pmatrix}
    \begin{pmatrix}
        f\\
        g 
    \end{pmatrix}
\end{equation*}
where $\Delta(x) = \theta(x)\Delta$ and $V(x) = H\delta(x)$.
\begin{enumerate}
    \item Prove $u_0^2 = 1 - v_0^2 = \frac{1}{2}(1 + \sqrt{E^2-\Delta^2}/E)$ and $\hbar k^{\pm} = \sqrt{2m(\mu \pm \sqrt{E^2 - \Delta^2})}$.
    \item Solve the scattering coefficient of this scattering problem.
\end{enumerate}
Hint: use plane-wave ansatz:
\begin{equation*}
    \begin{pmatrix}
        f\\
        g   
    \end{pmatrix}
    =
    \begin{pmatrix}
        u\\
        v
    \end{pmatrix}
    \eexp^{\ii (kx-Et/\hbar)}
\end{equation*}
then solve the stationary eigenvalue equation and apply the boundary condition at $x=0$.
\section*{Execise 3.3}
Prove the current conservation of quasi-particles in Andreev reflection problem:
\begin{equation*}
    \partial_t \rho + \nabla \cdot \vec{J}_\rho = 0
\end{equation*}
where $\rho=\abs{f}^2 + \abs{g}^2$ and $\vec{J}_\rho = \frac{\hbar}{m} \left[ \Im(f^*\nabla f) - \Im (g^*\nabla g) \right]$
\end{document}